\section{And One Thing That Shouldn't Be Possible}
\label{sec:in-browser}

Because the Go server compiles to WebAssembly and Postgres
(\href{https://pglite.dev}{pglite}) does too, Emi can run an
\textbf{entire backend inside a browser tab} --- the \emph{same} Go
handlers, the \emph{same} SQL, the \emph{same} generated DTO classes ---
talking to a real Postgres database, with no server process and no
network. It is not a separate, browser-only rewrite: compile the
Emi-generated Go server to \texttt{GOOS=js GOARCH=wasm}, give it an HTTP
router that lives in memory, and back its database calls with a real
Postgres compiled to WebAssembly. The browser then ``fetches'' from a
server sitting in the same tab.

\begin{sidebar}{In one sentence}
\texttt{main.js} calls \texttt{localFetch()} into Go WASM's
\texttt{net/http} router, which runs the ordinary generated handlers,
which query a \texttt{Database} interface backed by
\texttt{database-bridge.js} talking to pglite --- everything from the
handler down to the SQL string is byte-for-byte identical to what runs
on the real Gin server in \texttt{cmd/server}.
\end{sidebar}

\subsection{What is shared and what is swapped}
The whole design rests on one small interface --- handlers never import
\texttt{pgx} directly or hold a concrete connection, only a
\texttt{Database}:
\begin{lstlisting}
type Database interface {
  QueryRow(ctx context.Context, sql string, args ...any) pgx.Row
  Query(ctx context.Context, sql string, args ...any) (pgx.Rows, error)
  Exec(ctx context.Context, sql string, args ...any) (pgconn.CommandTag, error)
  CopyFrom(ctx context.Context, tableName pgx.Identifier,
    columnNames []string, rowSrc pgx.CopyFromSource) (int64, error)
}
\end{lstlisting}
Because this is exactly the shape of \texttt{*pgx.Conn}, the
\textbf{server build} just passes a live connection straight through
(passthrough). The \textbf{browser build} passes a \texttt{WasmDatabase}
that routes the same calls through a JS bridge into pglite instead ---
the handler code above this interface never knows, or needs to know,
which one it's talking to.

\begin{pullquote}
See it running: the
\href{https://torabian.github.io/emi/in-browser-server/}{in-browser
server demo} --- arguably the clearest possible demonstration of what
``one definition, real code everywhere'' actually buys a team: the exact
same handlers ship to production and run, unmodified, in a visitor's
browser tab.
\end{pullquote}
